\section{A triangle around each frozen intercept}
\label{sec:freezing}

We prove Theorem~\ref{thm:radial}. The required estimate is local to
one interval, but its error must be uniform over the whole partition.
Set
\[
 M=\max_J|\lambda|,\qquad A=LM,\qquad B=LM/2.
\]
Take any finite partition $\alpha=t_0<t_1<\cdots<t_N=\beta$.
For each interval $I=[t_{j-1},t_j]$, put $d_I=|I|$, let $m_I$ be
its midpoint, and set $R_I=\max_I R$. Write
$\delta=\max_I d_I$ for the mesh. The freezing tag $c_I\in I$
need not be the midpoint.

Move a point of \eqref{eq:radial-sweep} with $t\in I$ parallel to
$e+w(c_I)$ until it reaches $P$. Its intercept is
\[
 p_0+\rho(a+tb)+s(w(t)-w(c_I))=p_0+u a+v b,
\]
where, using the derivative hypothesis,
\begin{align}
 u&=\rho+s\int_{c_I}^t\lambda(\xi)\,d\xi,
 \label{eq:radial-coordinate}\\
 v-tu&=s\int_{c_I}^t\lambda(\xi)(\xi-t)\,d\xi.
 \label{eq:transverse-error}
\end{align}
Consequently
\begin{equation}\label{eq:old-footprint}
 -Ad_I\le u\le R_I+Ad_I,\qquad
 |v-m_Iu|\le\frac{d_I}{2}|u|+Bd_I^2.
\end{equation}
The first inequality uses $|t-c_I|\le d_I$; the second also uses
$|t-m_I|\le d_I/2$. In particular, the error transverse to a ray
is of order $d_I^2$, even though $u$ can move by order $d_I$ and can
become negative.

To enclose both the transverse error and the negative radial collar,
choose $C=2(A+B)+1$. Define the triangle in $P$
\begin{equation}\label{eq:triangle}
 T_I=\left\{p_0+u a+v b:
 -Cd_I\le u\le R_I+Ad_I,\quad
 |v-m_Iu|\le\frac{d_I}{2}(u+Cd_I)\right\}.
\end{equation}
Its width is zero at $u=-Cd_I$ and is positive at $u=R_I+Ad_I$.
Since $R_I+(A+C)d_I>0$, it is nondegenerate, including when $R_I=0$.

The triangle contains every intercept in
\eqref{eq:old-footprint}. The radial ranges are nested because
$C>A$. For $u\ge0$ the difference between the new half-width and
$d_I|u|/2$ is $Cd_I^2/2\ge Bd_I^2$. For $-Ad_I\le u<0$, that
difference is
\[
 d_Iu+\frac C2d_I^2\ge\left(\frac C2-A\right)d_I^2\ge Bd_I^2.
\]
This proves containment on both sides of the original radial tip;
Figure~\ref{fig:triangle} illustrates that extension.

\begin{figure}[htbp]
\centering
\begin{tikzpicture}[scale=1.3]
 \fill[blue!9] (-.675,-1.5)--(-.21,1.6)--(1.65,1.6)--cycle;
 \draw[blue!65!black,thick] (-.675,-1.5)--(-.21,1.6)--(1.65,1.6)--cycle;
 \fill[orange!30] (-.315,-.3)--(-.09,0)--(.15,1.6)--(1.29,1.6)--(.09,0)--(.045,-.3)--cycle;
 \draw[orange!70!black] (-.315,-.3)--(-.09,0)--(.15,1.6)--(1.29,1.6)--(.09,0)--(.045,-.3)--cycle;
 \draw[dashed] (-.675,-1.5)--(.72,1.6);
 \draw[-{Latex}] (-1,0)--(1.95,0) node[right] {$v$};
 \draw[-{Latex}] (0,-1.7)--(0,1.93) node[above] {$u$};
 \node[left] at (-.68,-1.45) {$-Cd_I$};
 \node[left] at (-.34,-.3) {$-Ad_I$};
 \node[right] at (1.65,1.6) {$R_I+Ad_I$};
 \node[blue!65!black] at (1.1,.45) {$T_I$};
\end{tikzpicture}
\caption{In intercept coordinates, the orange polygon bounds the
frozen sweep, including the small region below $u=0$. The enclosing
blue triangle extends its tip to $u=-Cd_I$; the dashed line is
$v=m_Iu$. The drawing is schematic; the triangle and footprint are
defined by \eqref{eq:triangle} and \eqref{eq:old-footprint}.}
\label{fig:triangle}
\end{figure}


Let $v_I=e+w(c_I)$ and take the base
$B_I=\pi_{v_I^\perp}T_I$. The restriction of this projection to
$P$ is nonsingular: its unit normal component is
$|e\cdot v_I|/\norm{v_I}=1/\norm{v_I}>0$.
Thus $B_I$ is again a nondegenerate compact triangle.
Moreover, $T_I+\R v_I=B_I+\R v_I$, since a point and its
orthogonal projection differ by a multiple of $v_I$.
The containment already proved shows that these cylinders cover
$\mathcal S$.

The area factor of projection from $P$ is $1/\norm{v_I}$.
Indeed, an orthonormal basis of $P$ projects to a parallelogram
whose area is the absolute scalar product of $e$ and
$v_I/\norm{v_I}$. The coordinate change $(u,v)\mapsto u a+v b$
has area factor $\kappa$. Integrating the width of
\eqref{eq:triangle} therefore gives the exact formula
\begin{equation}\label{eq:triangle-area}
 \area(B_I)=
 \frac{\kappa d_I\,[R_I+(A+C)d_I]^2}
      {2\sqrt{1+\norm{w(c_I)}^2}}.
\end{equation}
The collar changes $d_IR_I^2/2$ by $O(d_I^2)$, uniformly over $I$,
because $R$ is bounded. Since
$\sum_I d_I^2\le\delta(\beta-\alpha)$, its total contribution tends
to zero with the mesh. Put
$g(t)=(1+\norm{w(t)}^2)^{-1/2}$. Uniform continuity of $R$ and $g$
then gives
\[
 \sum_I\frac{\kappa d_I R_I^2g(c_I)}2
 \longrightarrow\frac\kappa2\int_J R(t)^2g(t)\,dt=\mathcal A.
\]
Thus the triangular covers have the asserted limit for arbitrary
partitions and tags. No convergence rate for the entire sum is needed
or asserted for an arbitrary continuous $R$.

For the actual projected bases, let $\mathcal S_I$ be the compact
sweep with $t\in I$, and put $B_I^{\mathrm{act}}=\pi_{v_I^\perp}\mathcal S_I$.
Moving parallel to $v_I$ does not change this projection, so the
intercept containment already proved gives $B_I^{\mathrm{act}}\subseteq B_I$.
The set $B_I^{\mathrm{act}}$ is compact and its cylinder contains
$\mathcal S_I$.

A matching lower bound comes from the sector at $s=0$:
\[
 E_I=\{p_0+\rho(a+tb):t\in I,\ 0\le\rho\le R(t)\}
 \subseteq\mathcal S_I.
\]
In the $(a,b)$ coordinates this map is $(\rho,t)\mapsto(\rho,\rho t)$,
which is one-to-one for $\rho>0$ and has absolute Jacobian $\rho$.
The edge $\rho=0$ maps to a single point. Consequently
$\area(E_I)=\frac\kappa2\int_I R(t)^2\,dt$, and projection from $P$
scales this area by $g(c_I)$. We obtain the two-sided estimate
\begin{equation}\label{eq:actual-sweep-squeeze}
 \frac{\kappa g(c_I)}2\int_I R(t)^2\,dt
 \le\area(B_I^{\mathrm{act}})\le\area(B_I).
\end{equation}
The sum of the lower bounds tends to $\mathcal A$: its difference
from the integral is at most
\[
 \frac\kappa2\sup_{\substack{r,t\in J\\|r-t|\le\delta}}
 |g(r)-g(t)|\int_J R(t)^2\,dt,
\]
which tends to zero by uniform continuity. The upper sums have the
same limit, proving the actual-base assertion, including when $R$
vanishes. This shared limit does not identify the individual actual
bases with their containing triangles or make their finite areas equal.
The proof of Theorem~\ref{thm:radial} is complete.
